% Volume IX: Institutional Design
% Part XVII. Designing for Delayed Truth
% Chapter 102: Rights as Computational Constraints
% Spine: proof-spines/ch102-spine.md

\chapter{Rights as Computational Constraints}
\label{ch:ch102}

Rights do not merely evaluate institutional action after optimization has finished; they limit the search space of permissible action from the start. If an institution can technically do many things, rights specify which of those things it is allowed even to consider as eligible outputs.

If \(\mathcal A\) is the set of technically possible actions, rights define
\[
\mathcal A_{\mathrm{legal}}
=
\{a\in\mathcal A:a\models\mathcal R\}.
\]
Optimization must therefore occur inside \(\mathcal A_{\mathrm{legal}}\), not across \(\mathcal A\) as a whole with rights added later as a moral gloss. On this view, rights are computational constraints on collection, inference, retention, and intervention---architectural limits on what the institution may treat as actionably available.
